% Part: first-order-logic
% Chapter: natural-deduction
% Section: proof-theoretic-notions

\documentclass[../../../include/open-logic-section]{subfiles}

\begin{document}

\iftag{FOL}
      {\olfileid[pt]{fol}{ntd}{ptn}}
      {\olfileid[pt]{pl}{ntd}{ptn}}

\olsection{Noções da Teoria da Prova}

\begin{editorial}
Esta seção reúne as definições da relação de demonstrabilidade
e de consistência para a dedução natural.
\end{editorial}

\begin{explain}
Assim como definimos várias noções semânticas importantes
(validade, consequência, satisfatibilidade), agora definimos noções
\emph{da Teoria da Prova} correspondentes.  Elas não são definidas por
apelo à satisfação de !!{sentence}s em !!{structure}s, mas por apelo à
!!{derivability} ou !!{nonderivability} de certas !!{sentence}s a
partir de outras.  Foi uma descoberta importante que essas noções
coincidem. Que coincidem é o conteúdo dos teoremas de
\emph{correção} e \emph{completude}.
\end{explain}


\begin{defn}[Teoremas]
Uma !!{sentence}~$!A$ é um \emph{teorema} se há uma !!{derivation}
de~$!A$ em dedução natural na qual todas as hipóteses são
!!{discharged}s.  Escrevemos $\Proves !A$ se $!A$ é um teorema e
$\Proves/ !A$ se não o é.
\end{defn}

\begin{defn}[!!^{derivability}]
!!^a{sentence} $!A$ é \emph{!!{derivable} de} um conjunto de
!!{sentence}s~$\Gamma$, $\Gamma \Proves !A$, se há uma
!!{derivation} com conclusão~$!A$ e na qual toda hipótese
é ou !!{discharged} ou está em~$\Gamma$. Se $!A$ não é
!!{derivable} de $\Gamma$, escrevemos $\Gamma \Proves/ !A$.
\end{defn}

\begin{defn}[Consistência]
Um conjunto de !!{sentence}s~$\Gamma$ é \emph{inconsistente} se, e somente se, $\Gamma
\Proves \lfalse$.  Se $\Gamma$ não é inconsistente, isto é, se
$\Gamma \Proves/ \lfalse$, dizemos que é \emph{consistente}.
\end{defn}

\begin{prop}[Reflexividade]
\ollabel{prop:reflexivity}
Se $!A \in \Gamma$, então $\Gamma \Proves !A$.
\end{prop}

\begin{proof}
A hipótese $!A$ por si só é uma !!{derivation} de~$!A$ em que toda
hipótese !!{undischarged} (isto é, $!A$) está em~$\Gamma$.
\end{proof}
  

\begin{prop}[Monotonicidade]
\ollabel{prop:monotonicity}
Se $\Gamma \subseteq \Delta$ e $\Gamma \Proves !A$, então $\Delta
\Proves !A$.
\end{prop}

\begin{proof}
Qualquer !!{derivation} de $!A$ a partir de $\Gamma$ é também uma !!{derivation} de
$!A$ a partir de~$\Delta$.
\end{proof}

\begin{prop}[Transitividade]
\ollabel{prop:transitivity}
Se $\Gamma \Proves !A$ e $\{!A\} \cup \Delta \Proves
!B$, então $\Gamma \cup \Delta \Proves !B$.
\end{prop}

\begin{proof}
Se $\Gamma \Proves !A$, há uma !!{derivation}~$\delta_0$ de~$!A$
com todas as hipóteses !!{undischarged}s em~$\Gamma$.  Se $\{!A\} \cup
\Delta \Proves !B$, então há uma !!{derivation}~$\delta_1$ de~$!B$
com todas as hipóteses !!{undischarged}s em~$\{!A\} \cup \Delta$.
Considere agora:
\begin{prooftree}
  \AxiomC{$\Delta, \Discharge{!A}{1}$}
  \RightLabel{$\delta_1$}
  \DeduceC{$!B$}
  \DischargeRule{\Intro{\lif}}{1}
  \UnaryInfC{$!A \lif !B$}
  \AxiomC{$\Gamma$}
  \RightLabel{$\delta_0$}
  \DeduceC{$!A$}
  \RightLabel{\Elim{\lif}}
  \BinaryInfC{$!B$}
\end{prooftree}
As hipóteses !!{undischarged}s estão agora todas entre $\Gamma \cup
\Delta$, de modo que isto mostra $\Gamma \cup \Delta \Proves !B$.
\end{proof}

Quando $\Gamma = \{!A_1, !A_2, \ldots, !A_k\}$ é um conjunto finito, podemos usar a notação simplificada $!A_1,!A_2,\ldots,!A_k \Proves !B$ para $\Gamma \Proves !B$; em particular, $!A \Proves !B$ significa que $\{!A\} \Proves !B$.

Note que se $\Gamma \Proves !A$ e $!A
\Proves !B$, então $\Gamma \Proves !B$. Segue também que se $!A_1,
\dots, !A_n \Proves !B$ e $\Gamma \Proves !A_i$ para cada~$i$, então
$\Gamma \Proves !B$.

\begin{prop}
\ollabel{prop:incons}
As seguintes afirmações são equivalentes.
    \begin{enumerate}
      \item \( \Gamma \) é inconsistente.
      \item \( \Gamma \Proves {!A} \) para toda !!{sentence}~\( {!A} \).
      \item \( \Gamma \Proves {!A} \) e \( \Gamma \Proves \lnot {!A} \) para alguma !!{sentence}~\( {!A} \).
    \end{enumerate}
\end{prop}

\begin{proof}
Exercício.
\end{proof}

\tagprob{FOL}
\begin{prob}
Prove \olref[fol][ntd][ptn]{prop:incons}
\end{prob}
\tagendprob

\tagprob{notFOL}
\begin{prob}
Prove \olref[pl][ntd][ptn]{prop:incons}
\end{prob}
\tagendprob

\begin{prop}[Compacidade]
  \ollabel{prop:proves-compact}
  \begin{enumerate}
  \item Se $\Gamma \Proves !A$, então há um subconjunto finito $\Gamma_0
    \subseteq \Gamma$ tal que $\Gamma_0 \Proves !A$.
  \item Se todo subconjunto finito de~$\Gamma$ é
    consistente, então $\Gamma$ é consistente.
  \end{enumerate}
\end{prop}

\begin{proof}
  \begin{enumerate}
    \item Se $\Gamma \Proves !A$, então há uma
      !!{derivation}~$\delta$ de~$!A$ a partir de~$\Gamma$. Seja $\Gamma_0$
      o conjunto de hipóteses !!{undischarged}s de~$\delta$.  Como
      toda !!{derivation} é finita, $\Gamma_0$ só pode conter
      finitamente muitas !!{sentence}s.  Assim, $\delta$ é uma !!{derivation}
      de~$!A$ a partir de um $\Gamma_0$ finito $\subseteq \Gamma$.
    \item Isto é a contrapositiva de (1) no caso especial $!A
      \ident \lfalse$.
  \end{enumerate}
\end{proof}

\end{document}
